\documentclass[11pt]{article}
\usepackage{mystyle}

\title{Rosen --- \S 6.2: The Pigeonhole Principle\\ \large Selected Exercises}
\author{Alston}
\date{Semester 2, 2026}

\begin{document}
\maketitle

% ======================================================================
\begin{exercise}{1}
    Show that in any set of six classes, each meeting regularly once a
    week on a particular day of the week, there must be two that meet
    on the same day, assuming that no classes are held on weekends.
\end{exercise}

% ======================================================================
\begin{exercise}{5}
    Show that among any group of five (not necessarily consecutive)
    integers, there are two with the same remainder when divided by
    4.
\end{exercise}

% ======================================================================
\begin{exercise}{7}
    Let $n$ be a positive integer. Show that in any set of $n$
    consecutive integers there is exactly one divisible by $n$.
\end{exercise}

% ======================================================================
\begin{exercise}{8}
    Show that if $f$ is a function from $S$ to $T$, where $S$ and $T$
    are finite sets with $|S| > |T|$, then there are elements $s_1$
    and $s_2$ in $S$ such that $f(s_1) = f(s_2)$, or in other words,
    $f$ is not one-to-one.
\end{exercise}

% ======================================================================
\begin{exercise}{9}
    What is the minimum number of students, each of whom comes from
    one of the 50 states, who must be enrolled in a university to
    guarantee that there are at least 100 who come from the same
    state?
\end{exercise}

% ======================================================================
\begin{exercise}{11}
    Let $(x_i, y_i, z_i)$, $i = 1, 2, 3, 4, 5, 6, 7, 8, 9$, be a set
    of nine distinct points with integer coordinates in $xyz$ space.
    Show that the midpoint of at least one pair of these points has
    integer coordinates.
\end{exercise}

% ======================================================================
\begin{exercise}{21}
    Construct a sequence of 16 positive integers that has no
    increasing or decreasing subsequence of five terms.
\end{exercise}

% ======================================================================
\begin{exercise}{22}
    Show that if there are 101 people of different heights standing
    in a line, it is possible to find 11 people in the order they are
    standing in the line with heights that are either increasing or
    decreasing.
\end{exercise}

% ======================================================================
\begin{exercise}{24}
    Suppose that 21 girls and 21 boys enter a mathematics competition.
    Furthermore, suppose that each entrant solves at most six
    questions, and for every boy-girl pair, there is at least one
    question that they both solved. Show that there is a question
    that was solved by at least three girls and at least three boys.
\end{exercise}

% ======================================================================
\begin{exercise}{29}
    \textit{Recall: Ramsey numbers were discussed after Example 13 in
    Section 6.2. $R(m,n)$ is the smallest integer such that any
    two-coloring of the edges of the complete graph $K_{R(m,n)}$
    contains either a red $K_m$ or a blue $K_n$.}

    Show that if $n$ is an integer with $n \geq 2$, then the Ramsey
    number $R(2,n)$ equals $n$.
\end{exercise}

% ======================================================================
\begin{exercise}{40}
    Prove that at a party where there are at least two people, there
    are two people who know the same number of other people there.
\end{exercise}

% ======================================================================
\begin{exercise}{41}
    An arm wrestler is the champion for a period of 75 hours. (Here,
    by an hour, we mean a period starting from an exact hour, such as
    1 p.m., until the next hour.) The arm wrestler had at least one
    match an hour, but no more than 125 total matches. Show that
    there is a period of consecutive hours during which the arm
    wrestler had exactly 24 matches.
\end{exercise}

% ======================================================================
\begin{exercise}{43}
    Show that if $f$ is a function from $S$ to $T$, where $S$ and $T$
    are nonempty finite sets and $m = \lceil |S|/|T| \rceil$, then
    there are at least $m$ elements of $S$ mapped to the same value
    of $T$. That is, show that there are distinct elements
    $s_1, s_2, \dots, s_m$ of $S$ such that
    $f(s_1) = f(s_2) = \cdots = f(s_m)$.
\end{exercise}

\end{document}